% Part: sets-functions-relations
% Chapter: relations
% Section: relations-as-sets

\documentclass[../../../include/open-logic-section]{subfiles}

\begin{document}

\olfileid[tr]{sfr}{rel}{set}
\olsection{Küme Olarak Bağıntılar}

\begin{explain}
\olref[sfr][set][imp]{sec} kesiminde bazı önemli kümelerden söz etmiştik:
$\Nat$, $\Int$, $\Rat$, $\Real$. Bu kümelerden bazılarının elemanları arasındaki
ilginç bağıntılardan birkaçını kuşkusuz anımsayacaksınız. Örneğin bu kümelerin
her biri üzerinde tümüyle standart bir \emph{sıralama bağıntısı} vardır. Ayrıca
her nesnenin kendisiyle ve yalnızca kendisiyle kurduğu \emph{kendisiyle özdeş
olma} bağıntısı bulunur. İleride daha birçok ilginç bağıntıyla karşılaşacağız; mümkün olan
bağıntıların sayısı ise daha da fazladır. Ancak bunları gözden geçirmeden önce,
bağıntılara özel bir küme türü olarak bakabileceğimize dikkat çekerek
başlayacağız.

Bunun için \olref[sfr][set][pai]{sec} kesiminden iki şeyi anımsayın. İlk olarak
\emph{sıralı ikili} kavramını anımsayın: $a$ ve $b$ verildiğinde
~$\tuple{a, b}$ ikilisini oluşturabiliriz. Burada elemanların sırası gerçekten
önemlidir. Dolayısıyla $a \neq b$ ise $\tuple{a, b} \neq \tuple{b, a}$ olur.
(Bunu sırasız ikililerle, yani $2$ elemanlı kümelerle karşılaştırın; bunlarda
$\{a, b\}=\{b, a\}$ olur.) İkinci olarak \emph{Kartezyen çarpım} kavramını
anımsayın: $A$ ve $B$ kümelerse, $x \in A$ ve $y \in B$ olan bütün
$\tuple{x, y}$ ikililerinin kümesi~$A \times B$'yi oluşturabiliriz. Özellikle
$A^{2}= A \times A$, elemanları $A$'dan gelen bütün sıralı ikililerin kümesidir.

Şimdi bir küme üzerindeki belirli bir bağıntıyı ele alacağız: doğal sayılar
kümesi~$\Nat$ üzerindeki $<$ bağıntısı. $n<m$ olan bütün $\tuple{n, m}$ sayı
ikililerinin kümesini düşünün; yani
\[
  R=\Setabs{\tuple{n, m}}{n, m \in \Nat \text{ ve } n<m}.
\]
$n$'nin $m$'den küçük olmasıyla $\tuple{n, m}$ ikilisinin $R$'nin bir elemanı
olması arasında şu yakın bağlantı vardır:
\[
      n<m\text{ ancak ve ancak }\tuple{n, m} \in R.
\]
Gerçekten de hiçbir bilgi yitirmeden, $R$ kümesini $\Nat$ üzerindeki $<$
bağıntısının kendisi olarak \emph{ele alabiliriz}.

Aynı biçimde, sayılar arasındaki herhangi bir bağıntı için $\Nat^{2}$'nin bir
alt kümesini kurabiliriz. Tersine, herhangi bir sayı ikilileri kümesi
$S \subseteq \Nat^{2}$ verildiğinde, sayılar arasında buna karşılık gelen bir
bağıntı vardır: $n$ ile $m$ arasındaki bağıntı, ancak ve ancak
$\tuple{n, m} \in S$ olduğunda geçerlidir. Bu da aşağıdaki tanımı gerekçelendirir:
\end{explain}

\begin{defn}[İkili bağıntı]
Bir $A$ kümesi üzerindeki \emph{ikili bağıntı}, $A^{2}$'nin bir alt kümesidir.
$R \subseteq A^{2}$, $A$ üzerinde bir ikili bağıntı ve $x, y \in A$ ise,
$\tuple{x, y} \in R$ yerine kimi zaman $Rxy$ (veya $xRy$) yazarız.
\end{defn}

\begin{ex}
  \ollabel{relations}
Doğal sayı ikililerinin kümesi $\Nat^{2}$, iki boyutlu bir matris biçiminde
şöyle listelenebilir:
\[
  \begin{array}{ccccc}
  \mathbf{\tuple{ 0,0 }} & \tuple{ 0,1 } &
    \tuple{ 0,2 } & \tuple{ 0,3 } & \ldots\\
  \tuple{ 1,0 } & \mathbf{\tuple{ 1,1 }} &
    \tuple{ 1,2 } & \tuple{ 1,3 } & \ldots\\
  \tuple{ 2,0 } & \tuple{ 2,1 } &
    \mathbf{\tuple{ 2,2 }} & \tuple{ 2,3 } & \ldots\\
  \tuple{ 3,0 } & \tuple{ 3,1 } & \tuple{ 3,2 } &
    \mathbf{\tuple{ 3,3 }} & \ldots\\
  \vdots & \vdots & \vdots & \vdots & \mathbf{\ddots}
  \end{array}
\]
Köşegen üzerinde bulunan ikililerden oluşan $\Nat^2$ alt kümesi, yani
\[
  \{\tuple{0,0 }, \tuple{ 1,1 }, \tuple{ 2,2 }, \dots\},
\]
$\Nat$ üzerindeki \emph{özdeşlik bağıntısı} olduğu için burada köşegeni kalın
yazdık. (Özdeşlik bağıntısı sık kullanıldığından, herhangi bir $A$ kümesi için
$\Id{A}=\Setabs{\tuple{ x,x }}{x \in A}$ biçiminde tanımlayalım.) Köşegenin
üstünde kalan bütün ikililerin alt kümesi, yani
\[
  L = \{\tuple{ 0,1 },\tuple{ 0,2 },\ldots,\tuple{ 1,2 },
  \tuple{ 1,3 }, \dots, \tuple{ 2,3 }, \tuple{ 2,4 },\ldots\},
\]
\emph{küçüktür} bağıntısıdır; başka bir deyişle $Lnm$ ancak ve ancak $n<m$
olduğunda geçerlidir. Köşegenin altında kalan ikililerin alt kümesi, yani
\[
  G=\{ \tuple{ 1,0 },\tuple{ 2,0 },\tuple{
    2,1 }, \tuple{ 3,0 },\tuple{ 3,1 },\tuple{ 3,2 }, \dots\},
\]
\emph{büyüktür} bağıntısıdır; yani $Gnm$ ancak ve ancak $n>m$ olduğunda
geçerlidir. $L$ ile $\Id{\Nat}$'ın birleşimi---$K=L\cup\Id{\Nat}$ diyebileceğimiz
bağıntı---\emph{küçük veya eşittir} bağıntısıdır: $Knm$ ancak ve ancak
$n \le m$. Benzer biçimde $H=G\cup\Id{\Nat}$, \emph{büyük veya eşittir}
bağıntısıdır. $L$, $G$, $K$ ve $H$ bağıntıları, \emph{sıralamalar} denilen özel
bağıntı türleridir. $L$ ile $G$'nin şöyle bir özelliği vardır: hiçbir sayı
kendisiyle $L$ veya $G$ bağıntısını kurmaz;
başka bir deyişle her $n$ için ne $Lnn$ ne de $Gnn$ geçerlidir. Bu özelliği
taşıyan bağıntılara \emph{yansımasız}; aynı zamanda sıralama iseler \emph{sıkı
sıralama} denir.
\end{ex}

\begin{explain}
Sıralamalar ve özdeşlik önemli ve doğal bağıntılar olsa da tanımımıza göre,
ne kadar yapay ya da zorlama görünürse görünsün, $A^{2}$'nin \emph{herhangi}
bir alt kümesinin $A$ üzerinde bir bağıntı olduğu vurgulanmalıdır. Özellikle
$\emptyset$ her küme üzerinde bir bağıntıdır (hiçbir eleman ikilisinin kurmadığı
\emph{boş bağıntı}); $A^{2}$'nin kendisi de $A$ üzerinde, her eleman ikilisinin
kurduğu \emph{evrensel bağıntı} denilen bir bağıntıdır. Bunun yanı sıra
$E=\Setabs{\tuple{n, m}}{n>5 \text{ veya } m \times n \ge 34}$ gibi bir küme de
bağıntı sayılır.
\end{explain}

\begin{prob}
  $\Pow{\{a, b, c\}}$ kümesi üzerindeki $\subseteq$ bağıntısının elemanlarını
  listeleyin.
\end{prob}

\end{document}
